\documentclass[12pt,a4paper]{article}
\usepackage{ctex}
\usepackage{amsmath,amssymb,amsthm}
\usepackage{geometry}
\usepackage{tcolorbox}
\usepackage{booktabs}
\usepackage{hyperref}
\usepackage{xcolor}

\geometry{left=2.5cm,right=2.5cm,top=2.5cm,bottom=2.5cm}

\tcbuselibrary{theorems,skins,breakable}

\newtcbtheorem[number within=section]{theorem}{定理}%
{colback=blue!5,colframe=blue!75!black,fonttitle=\bfseries,breakable}{th}

\newtcbtheorem[number within=section]{definition}{定义}%
{colback=green!5,colframe=green!50!black,fonttitle=\bfseries,breakable}{def}

\newtcbtheorem[number within=section]{example}{例子}%
{colback=orange!5,colframe=orange!75!black,fonttitle=\bfseries,breakable}{ex}

\newtcolorbox{keypoint}{colback=red!5!white,colframe=red!75!black,title=\textbf{核心要点},breakable}

\newtcolorbox{warning}{colback=yellow!5!white,colframe=yellow!75!black,title=\textbf{⚠️ 注意},breakable}

\newtcolorbox{summary}{colback=purple!5!white,colframe=purple!75!black,title=\textbf{总结},breakable}

\title{\textbf{可微性、恰当微分与可积条件}}
\author{关于状态函数与特性函数的深入讨论}
\date{\today}

\begin{document}

\maketitle

\begin{abstract}
本文详细讨论了可微性、全微分、恰当微分形式、可积条件以及路径无关性之间的关系。澄清了常见的误解，特别是：(1) 可微不要求偏导数连续；(2) 可积条件的充要性依赖于前提条件；(3) 状态函数与特性函数的区别。最后应用于分析力学中的拉格朗日量和哈密顿量，说明为什么需要勒让德变换。
\end{abstract}

\tableofcontents
\newpage

\section{核心概念的严格定义}

\subsection{可微性}

\begin{definition}{可微}{differentiable}
函数 $f(x, y)$ 在点 $(x_0, y_0)$ \textbf{可微}，如果存在常数 $A, B$ 使得：
\begin{equation}
f(x_0 + \Delta x, y_0 + \Delta y) - f(x_0, y_0) = A\Delta x + B\Delta y + o(\rho)
\end{equation}
其中 $\rho = \sqrt{(\Delta x)^2 + (\Delta y)^2}$，$o(\rho)$ 表示当 $\rho \to 0$ 时比 $\rho$ 高阶的无穷小。

如果 $f$ 可微，则：
\begin{equation}
A = \frac{\partial f}{\partial x}\bigg|_{(x_0, y_0)}, \quad B = \frac{\partial f}{\partial y}\bigg|_{(x_0, y_0)}
\end{equation}
\end{definition}

\begin{warning}
可微性\textbf{只要求}偏导数存在，\textbf{不要求}偏导数连续！
\end{warning}

\subsection{全微分}

\begin{definition}{全微分}{total-differential}
若 $f(x, y)$ 可微，则其\textbf{全微分}为：
\begin{equation}
df = \frac{\partial f}{\partial x}dx + \frac{\partial f}{\partial y}dy
\end{equation}
\end{definition}

\begin{keypoint}
任何可微函数都有全微分（这是定义）。
\end{keypoint}

\subsection{恰当微分形式}

\begin{definition}{恰当微分形式}{exact-differential}
微分形式 $\omega = P(x,y)\,dx + Q(x,y)\,dy$ 是\textbf{恰当的}（exact），如果存在函数 $f$ 使得：
\begin{equation}
\omega = df
\end{equation}
即：
\begin{equation}
P = \frac{\partial f}{\partial x}, \quad Q = \frac{\partial f}{\partial y}
\end{equation}
\end{definition}

\subsection{可积条件}

\begin{definition}{可积条件}{integrability-condition}
对于微分形式 $\omega = P\,dx + Q\,dy$，\textbf{可积条件}（也称恰当性条件、闭形式条件）为：
\begin{equation}
\boxed{\frac{\partial P}{\partial y} = \frac{\partial Q}{\partial x}}
\end{equation}
\end{definition}

\subsection{路径无关}

\begin{definition}{路径无关}{path-independence}
微分形式 $\omega$ 的积分\textbf{路径无关}，如果：
\begin{equation}
\int_A^B \omega
\end{equation}
只依赖于端点 $A, B$，不依赖于连接它们的路径。

等价地：
\begin{equation}
\oint_\gamma \omega = 0, \quad \forall \gamma \text{ 闭合}
\end{equation}
\end{definition}

\subsection{状态函数}

\begin{definition}{状态函数}{state-function}
函数 $F(x_1, \ldots, x_n)$ 是\textbf{状态函数}，如果满足：
\begin{enumerate}
\item \textbf{良定义}：每个状态点有唯一函数值
\item \textbf{全微分存在}：$dF = \sum_i \frac{\partial F}{\partial x_i}dx_i$
\item \textbf{可积条件}：$\frac{\partial^2 F}{\partial x_i \partial x_j} = \frac{\partial^2 F}{\partial x_j \partial x_i}$（若 $F \in C^2$）
\item \textbf{路径无关}：$\oint dF = 0$（在单连通域）
\item \textbf{瞬时性}：值只依赖当前状态，不依赖历史
\end{enumerate}
\end{definition}

\begin{keypoint}
核心等价关系：
\begin{equation}
\boxed{\text{状态函数} \Leftrightarrow \text{全微分存在} \Leftrightarrow \text{路径无关（单连通域）}}
\end{equation}
\end{keypoint}

\section{关键定理}

\subsection{可微保证全微分}

\begin{theorem}{可微$\Rightarrow$全微分}{diff-implies-total}
若 $f$ 可微，则 $df$ 存在且是恰当微分。
\end{theorem}

\begin{proof}
由定义，$df = \frac{\partial f}{\partial x}dx + \frac{\partial f}{\partial y}dy$ 就是 $f$ 的微分。
\end{proof}

\begin{warning}
此定理\textbf{不要求}偏导数连续！
\end{warning}

\subsection{可微保证路径无关}

\begin{theorem}{微积分基本定理}{fundamental-theorem}
若 $f$ 在单连通区域 $D$ 上可微，则：
\begin{equation}
\int_\gamma df = f(B) - f(A)
\end{equation}
路径无关。
\end{theorem}

\begin{proof}
这是微积分基本定理的推广。
\end{proof}

\begin{keypoint}
要求：
\begin{itemize}
\item ✅ $f$ 可微
\item ✅ $D$ 单连通
\item ❌ \textbf{不要求}偏导数连续
\end{itemize}
\end{keypoint}

\subsection{偏导数连续保证可微}

\begin{theorem}{充分条件}{sufficient-condition}
若 $\frac{\partial f}{\partial x}$ 和 $\frac{\partial f}{\partial y}$ 在 $(x_0, y_0)$ 的某邻域内存在，且在 $(x_0, y_0)$ 连续，则 $f$ 在 $(x_0, y_0)$ 可微。
\end{theorem}

\begin{warning}
这是\textbf{充分条件}，不是必要条件！

可微 $\not\Rightarrow$ 偏导数连续
\end{warning}

\subsection{可积条件判别恰当性（同济版）}

\begin{theorem}{可积条件的充要性}{tongji-theorem}
设 $D$ 是\textbf{单连通区域}，$P, Q \in C^1(D)$，则以下等价：
\begin{enumerate}
\item $\omega = P\,dx + Q\,dy$ 恰当
\item $\oint_\gamma \omega = 0$ 对所有闭合路径成立
\item 积分路径无关
\item \textbf{可积条件}：$\frac{\partial P}{\partial y} = \frac{\partial Q}{\partial x}$
\end{enumerate}
\end{theorem}

\begin{keypoint}
\textbf{关键前提}：
\begin{itemize}
\item ✅ $D$ 单连通
\item ✅ $P, Q \in C^1$（连续可微）
\end{itemize}

\textbf{在此前提下，可积条件是充要的！}
\end{keypoint}

\subsection{施瓦茨定理}

\begin{theorem}{施瓦茨定理}{schwarz}
若 $f \in C^2$（二阶连续可微），则：
\begin{equation}
\frac{\partial^2 f}{\partial x \partial y} = \frac{\partial^2 f}{\partial y \partial x}
\end{equation}
\end{theorem}

\begin{keypoint}
\textbf{推论}：若 $f \in C^2$，则可积条件自动满足。
\end{keypoint}

\section{层次关系}

\subsection{函数的光滑性层次}

\begin{equation}
C^\infty \subset C^2 \subset C^1 \subset \text{可微} \subset \text{连续}
\end{equation}

其中：
\begin{itemize}
\item $C^0$：连续函数
\item $C^1$：一阶偏导数存在且连续
\item $C^2$：一阶和二阶偏导数都存在且连续
\item $C^\infty$：所有阶偏导数都存在且连续（光滑）
\item 可微：偏导数存在（但可能不连续）
\end{itemize}

\subsection{概念的包含关系}

\begin{equation}
\text{函数} \supset \text{状态函数（全微分存在）} \supset \text{特性函数（偏导数给出动力学）}
\end{equation}

\subsection{条件的强弱}

\begin{table}[h]
\centering
\begin{tabular}{cccc}
\toprule
条件 & 保证可微 & 保证恰当 & 保证路径无关 \\
\midrule
$f$ 连续 & ✗ & ✗ & ✗ \\
$f$ 可微 & ✓（定义） & ✓ & ✓（单连通） \\
$f \in C^1$ & ✓ & ✓ & ✓（单连通） \\
$f \in C^2$ & ✓ & ✓ & ✓（单连通） \\
\bottomrule
\end{tabular}
\end{table}

\section{核心问题解答}

\subsection{可微是否保证偏导数连续？}

\textbf{答}：❌ \textbf{不保证}

\begin{example}{可微但导数不连续}{diff-not-continuous}
\begin{equation}
f(x) = \begin{cases}
x^2 \sin\frac{1}{x}, & x \neq 0 \\
0, & x = 0
\end{cases}
\end{equation}

\textbf{验证}：
\begin{itemize}
\item 可微：$f'(0) = \lim_{h \to 0} \frac{h^2\sin\frac{1}{h}}{h} = \lim_{h \to 0} h\sin\frac{1}{h} = 0$ ✓
\item 在 $x \neq 0$：$f'(x) = 2x\sin\frac{1}{x} - \cos\frac{1}{x}$
\item 当 $x \to 0$ 时，$\cos\frac{1}{x}$ 振荡，$\lim_{x \to 0} f'(x)$ 不存在
\item 所以 $f'$ 在 $x=0$ 不连续 ✗
\end{itemize}
\end{example}

\subsection{可微是否保证全微分存在？}

\textbf{答}：✅ \textbf{保证}

这是可微的\textbf{定义}！

\begin{equation}
\boxed{f \text{ 可微} \Rightarrow df = \frac{\partial f}{\partial x}dx + \frac{\partial f}{\partial y}dy \text{ 存在}}
\end{equation}

\subsection{可微是否保证恰当微分？}

\textbf{答}：✅ \textbf{保证}

\begin{equation}
\boxed{f \text{ 可微} \Rightarrow df \text{ 是恰当的}}
\end{equation}

因为 $df$ 就是 $f$ 的微分（定义）。

\subsection{可微是否保证路径无关？}

\textbf{答}：⚠️ \textbf{需要额外条件}

\begin{equation}
\boxed{f \text{ 可微} + D \text{ 单连通} \Rightarrow \int df \text{ 路径无关}}
\end{equation}

\textbf{不需要}偏导数连续。

\subsection{可积条件是充要条件吗？}

\textbf{答}：⚠️ \textbf{取决于前提}

\textbf{在标准前提下}（$P, Q \in C^1$，$D$ 单连通）：

\begin{equation}
\boxed{\frac{\partial P}{\partial y} = \frac{\partial Q}{\partial x} \Leftrightarrow \omega \text{ 恰当}}
\end{equation}

\textbf{是充要条件} ✓

\textbf{如果缺少前提}：
\begin{itemize}
\item 缺少 $C^1$：可积条件可能无法验证
\item 缺少单连通：可积条件只是必要条件，不充分
\end{itemize}

\subsection{``无法验证可积条件''等于``不恰当''吗？}

\textbf{答}：❌ \textbf{不等于}

\begin{equation}
\boxed{\text{``无法验证可积条件''} \neq \text{``不恰当''}}
\end{equation}

\begin{example}{恰当但无法验证}{exact-but-not-verifiable}
考虑 $f(x) = x^2\sin\frac{1}{x}$（$x \neq 0$），$f(0)=0$

\begin{itemize}
\item $df$ 恰当 ✓（因为 $f$ 存在）
\item 但 $f'$ 不连续，无法用 $C^1$ 框架下的可积条件验证
\end{itemize}
\end{example}

\section{物理应用：分析力学}

\subsection{拉格朗日量 $L(q, \dot{q}, t)$}

\textbf{状态空间}：切丛 $TQ = \{(q, \dot{q})\}$

\textbf{微分}：
\begin{equation}
dL = \frac{\partial L}{\partial q}dq + \frac{\partial L}{\partial \dot{q}}d\dot{q}
\end{equation}

\textbf{标准假设}：$L \in C^2$

\textbf{验证}：
\begin{itemize}
\item 可微 ✓
\item $dL$ 是恰当微分 ✓
\item 可积条件（施瓦茨定理）：$\frac{\partial^2 L}{\partial q \partial \dot{q}} = \frac{\partial^2 L}{\partial \dot{q} \partial q}$ ✓
\item 路径无关 ✓
\item \textbf{是切丛上的状态函数} ✓
\end{itemize}

\textbf{特性函数地位}：
\begin{itemize}
\item $L$ 是切丛的特性函数
\item 偏导数给出欧拉-拉格朗日方程：
\begin{equation}
\frac{d}{dt}\frac{\partial L}{\partial \dot{q}} - \frac{\partial L}{\partial q} = 0
\end{equation}
\end{itemize}

\subsection{哈密顿量 $H(q, p, t)$}

\textbf{状态空间}：相空间 $T^*Q = \{(q, p)\}$

\textbf{微分}：
\begin{equation}
dH = \frac{\partial H}{\partial q}dq + \frac{\partial H}{\partial p}dp
\end{equation}

\textbf{标准假设}：$H \in C^2$

\textbf{验证}：
\begin{itemize}
\item 可微 ✓
\item $dH$ 是恰当微分 ✓
\item 可积条件：$\frac{\partial^2 H}{\partial q \partial p} = \frac{\partial^2 H}{\partial p \partial q}$ ✓
\item 路径无关 ✓
\item \textbf{是相空间的状态函数} ✓
\end{itemize}

\textbf{特性函数地位}：
\begin{itemize}
\item $H$ 是相空间的特性函数
\item 偏导数给出正则方程：
\begin{equation}
\dot{q} = \frac{\partial H}{\partial p}, \quad \dot{p} = -\frac{\partial H}{\partial q}
\end{equation}
\end{itemize}

\subsection{代换后的 $L(q, p)$}

从 $p = \frac{\partial L}{\partial \dot{q}}$ 解出 $\dot{q}(q, p)$，代入：
\begin{equation}
L(q, p) = L(q, \dot{q}(q, p))
\end{equation}

\textbf{微分}：
\begin{equation}
dL = \frac{\partial L}{\partial q}dq + \frac{\partial L}{\partial p}dp
\end{equation}

\textbf{验证}：
\begin{itemize}
\item 可微 ✓（若原 $L \in C^2$ 且非退化）
\item $dL$ 是恰当微分 ✓
\item 可积条件满足 ✓
\item 路径无关 ✓
\item \textbf{是相空间的状态函数} ✓
\end{itemize}

\textbf{特性函数地位}：
\begin{itemize}
\item ❌ \textbf{不是}相空间的特性函数
\item 偏导数\textbf{不}给出正确的正则方程
\item $-\frac{\partial L}{\partial q} \neq \dot{p}$（差一个符号）
\end{itemize}

\subsection{三者对比}

\begin{table}[h]
\centering
\small
\begin{tabular}{lccc}
\toprule
\textbf{性质} & $L(q, \dot{q})$ & $L(q, p)$ & $H(q, p)$ \\
\midrule
空间 & 切丛 $TQ$ & 相空间 $T^*Q$ & 相空间 $T^*Q$ \\
可微 & ✓ & ✓ & ✓ \\
全微分存在 & ✓ & ✓ & ✓ \\
是状态函数 & ✓ & ✓ & ✓ \\
路径无关 & ✓ & ✓ & ✓ \\
可积条件 & ✓ & ✓ & ✓ \\
偏导数给出动力学 & ✓（E-L） & ✗ & ✓（正则） \\
是特性函数 & ✓（切丛） & ✗ & ✓（相空间） \\
物理意义 & 动能-势能 & 无明确意义 & 总能量 \\
\bottomrule
\end{tabular}
\end{table}

\section{为什么需要勒让德变换？}

\subsection{简单代换不够}

从 $L(q, \dot{q})$ 到 $L(q, p)$：只是变量代换

\begin{itemize}
\item 得到状态函数 ✓
\item 但不是特性函数 ✗
\end{itemize}

\subsection{勒让德变换的作用}

\begin{equation}
H(q, p) = p\dot{q} - L(q, \dot{q})|_{\dot{q}(p)}
\end{equation}

\textbf{不只是代换}，还有额外项 $p\dot{q}$！

\textbf{效果}：
\begin{itemize}
\item 从 $L = T - V$ 变成 $H = T + V$
\item 符号改变使得偏导数满足正则方程
\end{itemize}

\textbf{验证}：
\begin{align}
\frac{\partial H}{\partial p} &= \dot{q} + p\frac{\partial \dot{q}}{\partial p} - \frac{\partial L}{\partial \dot{q}}\frac{\partial \dot{q}}{\partial p} = \dot{q}
\end{align}

（因为 $p = \frac{\partial L}{\partial \dot{q}}$，两项抵消）

\subsection{总结}

\begin{summary}
\begin{equation}
\boxed{
\begin{array}{c}
\text{简单代换：} L(q,\dot{q}) \to L(q,p) \\
\Downarrow \\
\text{得到状态函数，但非特性函数} \\
\
\text{勒让德变换：} L(q,\dot{q}) \to H(q,p) \\
\Downarrow \\
\text{得到特性函数，偏导数给出动力学}
\end{array}
}
\end{equation}
\end{summary}

\section{关键结论总结}

\subsection{关于可微与连续性}

\begin{equation}
\boxed{
\begin{array}{rcl}
\text{偏导数连续} &\Rightarrow& \text{可微（充分）} \\
\text{可微} &\not\Rightarrow& \text{偏导数连续}
\end{array}
}
\end{equation}

\subsection{关于恰当性}

\begin{equation}
\boxed{
\begin{array}{rcl}
f \text{ 可微} &\Rightarrow& df \text{ 恰当（定义）} \\
\text{偏导数不连续} &\not\Rightarrow& df \text{ 不恰当}
\end{array}
}
\end{equation}

\subsection{关于路径无关}

\begin{equation}
\boxed{
\begin{array}{l}
f \text{ 可微 + 单连通} \Rightarrow \int df \text{ 路径无关} \\
\text{不需要偏导数连续}
\end{array}
}
\end{equation}

\subsection{关于可积条件}

\begin{equation}
\boxed{
\begin{array}{c}
P, Q \in C^1 + D \text{ 单连通} \\
\Downarrow \Updownarrow \\
\frac{\partial P}{\partial y} = \frac{\partial Q}{\partial x} \Leftrightarrow \omega \text{ 恰当} \\
\text{（充要条件）}
\end{array}
}
\end{equation}

\textbf{缺少前提时}：
\begin{itemize}
\item 可积条件可能无法验证（若 $P, Q \notin C^1$）
\item 或只是必要不充分（若非单连通）
\end{itemize}

\subsection{关于状态函数 vs 特性函数}

\begin{equation}
\boxed{
\begin{array}{c}
\text{状态函数：全微分存在，路径无关} \\
\Downarrow \text{（更强要求）} \\
\text{特性函数：偏导数给出动力学方程}
\end{array}
}
\end{equation}

\textbf{在相空间}：
\begin{itemize}
\item $L(q,p)$：是状态函数 ✓，非特性函数 ✗
\item $H(q,p)$：是状态函数 ✓，是特性函数 ✓
\end{itemize}

\section{常见误解澄清}

\subsection{误解1：可微需要偏导数连续}

✗ \textbf{错误}

✓ \textbf{正确}：可微只需要偏导数存在，不需要连续

\textbf{反例}：$f(x) = x^2\sin\frac{1}{x}$

\subsection{误解2：恰当微分需要偏导数连续}

✗ \textbf{错误}（如果已知原函数）

✓ \textbf{正确}：
\begin{itemize}
\item 若已知 $f$ 可微，$df$ 恰当（不需偏导数连续）
\item 若要\textbf{判断} $\omega$ 是否恰当，需要 $P, Q \in C^1$ 来验证可积条件
\end{itemize}

\subsection{误解3：可积条件总是充要的}

✗ \textbf{错误}

✓ \textbf{正确}：可积条件是充要的\textbf{需要前提}（$C^1$ + 单连通）

\subsection{误解4：无法验证可积条件就是不恰当}

✗ \textbf{错误}

✓ \textbf{正确}：无法验证 $
eq$ 不恰当

\subsection{误解5：$L(q,p)$ 不是状态函数}

✗ \textbf{错误}

✓ \textbf{正确}：$L(q,p)$ \textbf{是}状态函数，但\textbf{不是}特性函数

\section{完整逻辑链}

\begin{verbatim}
函数 f
  ↓ （定义）
可微？
  ├─ 是 → 偏导数存在
  │      ↓
  │    df 存在（全微分）
  │      ↓
  │    df 恰当（定义）
  │      ↓ （单连通）
  │    ∫df 路径无关
  │      ↓
  │    f 是状态函数
  │
  └─ 否 → df 不存在

若 f ∈ C¹（偏导数连续）
  ↓
可用可积条件验证恰当性

若 f ∈ C²
  ↓
可积条件自动满足（施瓦茨定理）
  ↓
是状态函数
  ↓ （若偏导数有物理意义）
可能是特性函数
\end{verbatim}

\section{核心要点总结}

\begin{keypoint}
\begin{enumerate}
\item 可微 $\neq$ 偏导数连续
\item 可微 $\Rightarrow$ 恰当（定义）
\item 可积条件充要性需要前提（$C^1$ + 单连通）
\item 状态函数 $\supset$ 特性函数
\item $L(q,p)$ 是状态函数但非特性函数
\item $H(q,p)$ 既是状态函数又是特性函数
\item 勒让德变换不只是变量代换，还改变函数形式
\end{enumerate}
\end{keypoint}

\end{document}